\clearpage
\VerdierPhysicalPageAnchor{0183}%
\begingroup\small\setlength{\parskip}{0pt}\setlength{\emergencystretch}{1.5em}
\setlength{\abovedisplayskip}{4pt plus 2pt minus 2pt}%
\setlength{\belowdisplayskip}{4pt plus 2pt minus 2pt}%

\VerdierSemanticAnchor{ch3d-1-3-2-morphism-m-cont}
Le morphisme $m$ est un monomorphisme et on a la suite exacte :

\VerdierSemanticAnchor{ch3d-1-3-2-exact-sequence}
\begin{equation}
0\longrightarrow c(\operatorname{id}_X)
\xrightarrow{m}c(\mu(S))[-1]
\longrightarrow c(\operatorname{id}_Z)[-1]
\longrightarrow0.
\tag{1.3.2.3}
\end{equation}
Les complexes $c(\operatorname{id}_X)$ et $c(\operatorname{id}_Z)$ sont
homotopes à zéro. Par suite, le complexe $c(\mu(S))$ est acyclique. Le
morphisme $\mu(S)$ est donc un quasi-isomorphisme.

\VerdierSemanticAnchor{ch3d-1-3-2-k-diagram}
En passant dans la catégorie $K(\mathcal A)$, on obtient un diagramme
(chap. I, 3.2.2.7) :
\begin{center}
\resizebox{0.84\textwidth}{!}{%
$\begin{tikzcd}[ampersand replacement=\&,column sep=3.2em,row sep=2.6em]
X \arrow[r,"\dot u"]
  \& Y \arrow[r,"\dot{p}(u)"] \arrow[dr,"\dot v"']
  \& c(\dot u) \arrow[r,"\dot{q}(u)"] \arrow[d,"\dot\mu(S)"]
  \& X[1] \\
\& \& Z \&
\end{tikzcd}$%
}
\end{center}
Posons alors :

\VerdierSemanticAnchor{ch3d-1-3-2-connecting}
\begin{equation}
\delta(S)=[q(u)][\mu(S)]^{-1}.
\tag{1.3.2.4}
\end{equation}
Il résulte des définitions des structures triangulées de $K(\mathcal A)$
(chap. I, 3.3.1) et de $D(\mathcal A)$ (chap. II, 2.2.6) que le triangle
$(X,Y,Z,[u],[v],\delta(S))$ est distingué dans $D(\mathcal A)$. De plus,
on vérifie immédiatement que le morphisme $\delta(S)$ est fonctoriel par
rapport à $S$ ((1.3.1), propriété a)). On a donc muni le foncteur canonique
$\operatorname{comp}(\mathcal A)\longrightarrow D(\mathcal A)$ d'une
structure de $\delta$-foncteur. Le $\delta$-foncteur ainsi obtenu est appelé
le \emph{$\delta$-foncteur canonique}. Le morphisme $\delta(S)$ est appelé
le \emph{morphisme de liaison} associé à la suite exacte $S$.

\VerdierSemanticAnchor{ch3d-1-3-3}
\medskip\noindent\textbf{1.3.3.}\quad
Soit $\mathcal A^o$ la catégorie opposée à $\mathcal A$. Pour toute suite
exacte $S$ de $\operatorname{comp}(\mathcal A^o)$, désignons par
$\delta^o(S)$ le morphisme de liaison associé. Soient $C$ une catégorie,
$X$ un objet de $C$ et $u$ un morphisme de $C$. Dans la proposition
ci-après, nous désignerons abusivement par $X^o$ et $u^o$ l'objet et le
morphisme correspondant de la catégorie opposée à $C$ (rappelons que la
catégorie opposée à $C$ a même ensemble d'objets et de morphismes que la
catégorie $C$).

\VerdierSemanticAnchor{ch3d-1-3-4}
\medskip\noindent\textbf{Proposition 1.3.4.}\quad
\textit{a) Pour toute suite exacte $S$ de
$\operatorname{comp}(\mathcal A)$, on a :}
\[
\delta(S[1])=-\delta(S)[1].
\]
\noindent\textit{b) Pour toute suite exacte $S$ de
$\operatorname{comp}(\mathcal A)$, compte tenu de l'identification
(1.2.8), on a :}
\[
\delta^o(S^o)=\delta(S)^o[1].
\]
\endgroup
